\documentclass[12pt,a4paper]{article}
\usepackage{amsmath,amscd,amsbsy,amssymb,latexsym,url,bm,amsthm}
\usepackage{epsfig,graphicx,subfigure}
\usepackage{enumitem,balance}
\usepackage{wrapfig}
\usepackage{mathrsfs,euscript}
\usepackage[usenames]{xcolor}
\usepackage{hyperref}
\usepackage[vlined,ruled,commentsnumbered,linesnumbered]{algorithm2e}

\newtheorem{theorem}{Theorem}
\newtheorem{lemma}[theorem]{Lemma}
\newtheorem{proposition}[theorem]{Proposition}
\newtheorem{corollary}[theorem]{Corollary}
\newtheorem{exercise}{Exercise}
\newtheorem*{solution}{Solution}
\newtheorem{definition}{Definition}
\theoremstyle{definition}


%\numberwithin{equation}{section}
%\numberwithin{figure}{section}

\renewcommand{\thefootnote}{\fnsymbol{footnote}}

\newcommand{\postscript}[2]
 {\setlength{\epsfxsize}{#2\hsize}
  \centerline{\epsfbox{#1}}}

\renewcommand{\baselinestretch}{1.0}

\setlength{\oddsidemargin}{-0.365in}
\setlength{\evensidemargin}{-0.365in}
\setlength{\topmargin}{-0.3in}
\setlength{\headheight}{0in}
\setlength{\headsep}{0in}
\setlength{\textheight}{10.1in}
\setlength{\textwidth}{7in}
\makeatletter \renewenvironment{proof}[1][Proof] {\par\pushQED{\qed}\normalfont\topsep6\p@\@plus6\p@\relax\trivlist\item[\hskip\labelsep\bfseries#1\@addpunct{.}]\ignorespaces}{\popQED\endtrivlist\@endpefalse} \makeatother
\makeatletter
\renewenvironment{solution}[1][Solution] {\par\pushQED{\qed}\normalfont\topsep6\p@\@plus6\p@\relax\trivlist\item[\hskip\labelsep\bfseries#1\@addpunct{.}]\ignorespaces}{\popQED\endtrivlist\@endpefalse} \makeatother



\begin{document}
\noindent

%========================================================================
\noindent\framebox[\linewidth]{\shortstack[c]{
\Large{\textbf{Exercise Sheet 18}}\vspace{1mm}\\
Discrete Mathematics by Hongfei Fu, 2023.12.07}}

\begin{center}
	\footnotesize{\color{black} \quad Name:\_\_\_\_\_\_\_\_\_  \quad Student ID:\_\_\_\_\_\_\_\_\_ \quad Class: \_\_\_\_\_\_\_\_\_\_\_\_} \\
	

\end{center}


\begin{enumerate}
\item  
In Exercises (3), (4) use Prim's algorithm to find a minimum spanning tree for the given weighted graph.\\
3. \includegraphics*[scale=0.8]{3@p802.png} \qquad \qquad 
4. \includegraphics*[scale=0.5]{4@p802.png}\\
\textbf{Answer Area:}
\vspace{6cm}

\item 
Use Kruskal's algorithm to find a minimum spanning tree for the weighted graph in Exercise(3), (4) respectively.

\textbf{Answer Area:}
\vspace{8cm}
\clearpage

\item 
Find a maximum spanning tree for the weighted graph in Exercise 3, 4 respectively.


\textbf{Answer Area:}
\vspace{7cm}

\item 
Show that there is a unique minimum spanning tree in a connected weighted graph if the weights of the edges are all different.


\textbf{Answer Area:}
\vspace{6cm}

\item (Additional Homework, \textbf{2pt})
Suppose $G = (V,E)$ a weighted connected simple graph with weight function $w: E \to \mathbb{R}$, and $e_0 \in E$ is an edge that has the greatest weight, i.e. for any $e \in E \setminus \{ e_0 \}$, $w(e_0) > w(e)$. Prove that if $e_0$ is in a minimum spanning tree of $G$, then $e_0$ must be a cut edge of $G$.

\textbf{Answer Area:}

\end{enumerate}
%========================================================================
\end{document}
